\chapter{Минимальный родительский функционал сдвига Штюкельберга}
\label{ch:v9-kms-minimal-stueckelberg}

\section{Удвоенный even carrier}

Чтобы получить nontrivial shift-invariant parent, введём одну compensator
copy десятимерного type/package carrier:
\begin{equation}
 X_{\rm St}=E_x\oplus E_y,
 \qquad
 E_x\simeq E_y\simeq V_{\rm type}\otimes P_{\rm KMS},
 \qquad \dim_{\mathbb C}X_{\rm St}=20.
 \label{eq:v9-kms-stueckelberg-carrier}
\end{equation}
Diagonal shift действует как
\begin{equation}
 x\longmapsto x+\varepsilon,
 \qquad
 y\longmapsto y+\varepsilon,
 \qquad \varepsilon\in E.
 \label{eq:v9-kms-stueckelberg-shift}
\end{equation}
Его generator-map равен
\begin{equation}
 T=\begin{pmatrix}I_{10}\\I_{10}\end{pmatrix},
 \qquad \rank T=10.
 \label{eq:v9-kms-stueckelberg-orbit-map}
\end{equation}

\section{Gauge-invariant quotient}

Минимальная invariant combination есть
\begin{equation}
 z=C\binom{x}{y}=x-y,
 \qquad
 C=\begin{pmatrix}I_{10}&-I_{10}\end{pmatrix},
 \qquad CT=0.
 \label{eq:v9-kms-stueckelberg-invariant}
\end{equation}
Поэтому nontrivial bounded parent задаётся
\begin{equation}
 S_{\rm St}(x,y)
 =\frac12(x-y)^TD_{\rm aux}(x-y),
 \qquad D_{\rm aux}=R_\theta\oplus R_\kappa>0.
 \label{eq:v9-kms-stueckelberg-parent}
\end{equation}
Его Hessian имеет block-форму
\begin{equation}
 H_{\rm St}=C^TD_{\rm aux}C
 =\begin{pmatrix}
 D_{\rm aux}&-D_{\rm aux}\\
 -D_{\rm aux}&D_{\rm aux}
 \end{pmatrix}.
 \label{eq:v9-kms-stueckelberg-hessian}
\end{equation}
Точные rank и nullity:
\begin{equation}
 \rank H_{\rm St}=10,
 \qquad
 \dim\ker H_{\rm St}=10.
 \label{eq:v9-kms-stueckelberg-rank-nullity}
\end{equation}
Причём
\begin{equation}
 H_{\rm St}T=0,
 \qquad
 \ker H_{\rm St}=\operatorname{im}T.
 \label{eq:v9-kms-stueckelberg-kernel-orbit}
\end{equation}
В isotropic point $D_{\rm aux}=I_{10}$ spectrum равен
\begin{equation}
 \Spec H_{\rm St}=\{0^{(10)},2^{(10)}\},
 \label{eq:v9-kms-stueckelberg-spectrum}
\end{equation}
поэтому parent positive semidefinite и не имеет лишних zero modes.

\section{Gauge fixing и FP determinant}

Выберем linear condition
\begin{equation}
 F(x,y)=D_{\rm aux}x.
 \label{eq:v9-kms-stueckelberg-condition}
\end{equation}
Её derivative вдоль orbit равна
\begin{equation}
 M_{\rm FP}
 =\frac{\partial F(x+\varepsilon,y+\varepsilon)}
 {\partial\varepsilon}\bigg|_{\varepsilon=0}
 =D_{\rm aux},
 \qquad \rank M_{\rm FP}=10.
 \label{eq:v9-kms-stueckelberg-fp}
\end{equation}
Следовательно,
\begin{equation}
 \det M_{\rm FP}
 =\det R_\theta\,\det R_\kappa.
 \label{eq:v9-kms-stueckelberg-fp-determinant}
\end{equation}
На уровне одной ghost sector это ровно target.

\section{Quotient obstruction}

Однако quotient не пуст:
\begin{equation}
 \dim_{\mathbb C}(X_{\rm St}/\operatorname{im}T)
 =20-10=10,
 \label{eq:v9-kms-stueckelberg-quotient-dimension}
\end{equation}
и представлен gauge-invariant boson $z=x-y$. Если он интегрируется как
complex Gaussian, то
\begin{equation}
 Z_z\propto(\det D_{\rm aux})^{-1}.
 \label{eq:v9-kms-stueckelberg-boson-factor}
\end{equation}
Ghost sector даёт
\begin{equation}
 Z_{\rm gh}=\det D_{\rm aux}.
 \label{eq:v9-kms-stueckelberg-ghost-factor}
\end{equation}
Полный determinant нейтрален:
\begin{equation}
 Z_zZ_{\rm gh}=1,
 \qquad
 -\log(Z_zZ_{\rm gh})=0.
 \label{eq:v9-kms-stueckelberg-cancellation}
\end{equation}
Если же $z$ не интегрировать, он остаётся десятимёрным новым
gauge-invariant sector:
\begin{equation}
 N_{\rm new\ quotient\ modes}=10.
 \label{eq:v9-kms-stueckelberg-new-modes}
\end{equation}
Таким образом, route имеет строгую развилку
\begin{equation}
 \text{integrate }z\Rightarrow\text{logdet cancellation},
 \qquad
 \text{retain }z\Rightarrow\text{new physical/auxiliary quotient}.
 \label{eq:v9-kms-stueckelberg-dichotomy}
\end{equation}

\section{ProofDSL и итог}

ProofDSL проверяет orbit/quotient ranks, exact invariance, Hessian
rank/nullity, zero-mode equality, positive isotropic spectrum, FP determinant
и exact determinant cancellation:
\begin{equation}
 N_{\rm ProofDSL\ obligations}=10/10,
 \qquad
 N_{\rm registered\ gates}=55,
 \qquad
 N_{\rm registered\ obligations}=432.
 \label{eq:v9-kms-stueckelberg-proofdsl}
\end{equation}
Итоговые ledgers:
\begin{equation}
 \begin{aligned}
 N_{\rm Stueckelberg\ gauge\ architecture}&=10/10,\\
 N_{\rm nontrivial\ rank\ ten\ orbit}&=1/1,\\
 N_{\rm no\ new\ quotient\ modes}&=0/1,\\
 N_{\rm uncancelled\ target\ logdet}&=0/1,\\
 N_{\rm physical\ four\ slot\ parent}&=0/1.
 \end{aligned}
 \label{eq:v9-kms-stueckelberg-final-ledger}
\end{equation}

\begin{theorem}[Минимальный nontrivial shift-parent]
Одна compensator copy даёт bounded Stückelberg parent с exact
rank-ten gauge orbit. Hessian имеет rank/nullity $10/10$, а его kernel
совпадает с orbit. FP determinant равен target determinant.
\end{theorem}

\begin{nogocriterion}[Stückelberg quotient отменяет логарифм определителя]
Та же минимальная конструкция оставляет десятимёрную gauge-invariant
combination $z=x-y$. Её complex bosonic determinant является обратным к
ghost determinant. Поэтому интегрирование отменяет target логарифм определителя, а
неинтегрированный $z$ добавляет новый quotient sector. Ни одна ветвь не
закрывает physical parent.
\end{nogocriterion}

\begin{protocol}\begin{enumerate}
\item even Stückelberg carrier: \textbf{$20$};
\item gauge orbit rank: \textbf{$10$};
\item Hessian rank/nullity: \textbf{$10/10$};
\item quotient dimension: \textbf{$10$};
\item FP determinant: \textbf{$\det R_\theta\det R_\kappa$};
\item combined determinant: \textbf{$1$};
\item ProofDSL: \textbf{$10/10$, registry $55/432$};
\item physical логарифм определителя closure: \textbf{$0/1$};
\item следующий гейт: \textbf{physical fermion-loop parent-origin}.
\end{enumerate}\end{protocol}

Машинный сертификат сохранён в
\path{s2t_v9_endpoint_creation_kms_logdet_minimal_stueckelberg_shift_parent_architecture_gate_results.json}.

\begin{thebibliography}{9}
\bibitem{v9-kms-stueckelberg-review}
H.~Ruegg, M.~Ruiz-Altaba,
\emph{The Stueckelberg field},
Int. J. Mod. Phys. A \textbf{19} (2004), 3265--3348,
arXiv:hep-th/0304245.
\bibitem{v9-kms-stueckelberg-brst}
N.~Boulanger, C.~Deffayet, S.~Garcia-Saenz, L.~Traina,
\emph{Consistent deformations of free massive field theories in the
Stueckelberg formulation},
JHEP \textbf{07} (2018) 021, arXiv:1806.04695.
\end{thebibliography}